\section{Results}

\subsection{Campaign validity}

The main campaign executed \NumRuns{} runs over \NumContexts{} contexts with
\NumPolicies{} policies and \NumSeeds{} seeds. No run failed. Every run
completed every vehicle it generated (minimum completion rate
\NumMinCompletion{}), and there were \NumTeleports{} teleports in the entire
campaign, so no result rests on a vehicle the simulator removed. All
\NumContextsAnalysed{} contexts satisfy the validity and completeness rules
declared in advance, and none is excluded.

\subsection{The portfolio is complementary, and one policy is near-dominant}

Gate G1$'$ --- declared before the campaign and reported here whichever way it
came out --- \NumGOneVerdict{}: \NumGOneNPolicies{} policies
(\NumGOnePolicies) are strictly best in at least one context with a resolved
margin. \textbf{P2 is never a resolved winner anywhere in the factor space.}
Resolved winner counts are P3 \NumWinResPThree{}, P4 \NumWinResPFour{} and P1
\NumWinResPOne{}.

The winner map is multi-factor but shallow. Always naming the same policy is
correct in \NumWinMapConst{} of contexts; the best single factor
(\NumWinMapOneFactor) reaches \NumWinMapOne{}, and the best pair of factors
\NumWinMapTwo{}. The regime structure is interpretable and matches the design:
P1 wins only where the shortest corridor is uncongested \emph{and} penetration
is high enough for the adaptive policies to over-divert; P4 wins at low demand
and generous green; P3 wins every context at the highest demand level.

\subsection{Adaptation is worth almost nothing --- but the portfolio is worth a
great deal}

These are different quantities and the difference is the central result.

The single best fixed policy is \NumGlobalSBS{}, optimal in \NumSBSOptimalPct{}
of contexts. Running it always costs \NumCostSBS~s of system mean journey time
per vehicle; the retrospective best-per-context choice costs \NumCostVBS~s. The
entire value of adapting the policy to the context is therefore
\textbf{\NumHeadroom~s per vehicle}, or \NumHeadroomPct{}, with a maximum of
\NumHeadroomMax~s in any single context.

Choosing the \emph{wrong fixed policy} costs far more. Always running the
distance-minimising policy costs \NumFixedPenPOne~s per vehicle more than
always running the load-balancing policy --- \NumFixedPenPctPOne{}. The
reactive and reliability-aware policies cost \NumFixedPenPTwo~s
(\NumFixedPenPctPTwo{}) and \NumFixedPenPFour~s (\NumFixedPenPctPFour{})
respectively.

\begin{quote}
Getting the fixed policy right is worth up to \NumFixedWorstPenaltyPct{} of
journey time. Getting the \emph{context-dependent} choice right, given that the
fixed policy is already the best one, is worth \NumHeadroomPct{}.
\end{quote}

\subsection{Why: an eighteen-fold asymmetry}

The explanation is not that the policies behave alike. It is that the best
fixed policy is near-dominant in a specific, measurable sense. In
\NumAsymNotBestPct{} of contexts \NumGlobalSBS{} is not the best policy, and
when it is not, it loses \NumAsymLoss~s on average (at most
\NumAsymMaxLoss~s). In \NumAsymBestPct{} of contexts it is the best, and when
it is, it wins by \NumAsymMargin~s on average (at most \NumAsymMaxMargin~s).
Its upside is \textbf{\NumAsymRatio$\times$ its downside}.

A portfolio containing such a policy has a large \emph{ranking} structure and a
small \emph{decision} structure. Any selector that moves away from it risks a
large loss to chase a small gain.

\subsection{The effect is at the edge of what the experiment can resolve}

Among the \NumResolvNContexts{} contexts where the fixed policy is not best, the
mean headroom is \NumResolvHeadroomNZ~s against a mean two-standard-error noise
band of \NumResolvNoiseNZ~s on the same paired comparison, and the headroom
exceeds its own noise in \NumResolvExceeds{} of them (\NumResolvExceedsPct{}).
Resolving the campaign-wide mean effect of \NumResolvHeadroom~s would require
roughly \NumResolvSeedsNeeded{} seeds per context --- about
\NumResolvRunsNeeded{} runs. We report this rather than presenting a
sub-noise-floor improvement as an achievement.

\subsection{The mechanistic rule, and why accuracy is not value}

Fitted on all contexts for exposition, the depth-2 mechanistic rule splits
first on $\NumMechFeature$, the share of the shortest corridor's capacity that
the guided cohort alone would consume, at a threshold bracketed between
\NumMechBracketLo{} and \NumMechBracketHi{} --- reported as a bracket, since
the grid cannot place it more finely. The rule is physically readable: when the
controlled flow is a small share of the arterial's capacity the decision turns
on the effective green ratio, and when it is a large share it turns on whether
the arterial is already oversaturated.

The rule identifies the best policy in \NumMechTopOne{} of contexts, more often
than always naming the fixed policy. Its mean decision regret is nevertheless
\NumMechRegret~s against \NumMechFixedRegret~s for the fixed policy.
\textbf{It is more accurate and more expensive.} Under leave-one-context-out
evaluation the same pattern holds: B1 reaches \NumBOneTopOne{} top-1 accuracy
against B0's \NumBZeroTopOne{}, and \NumBOneRegret~s of regret against B0's
\NumBZeroRegret~s.

This is the predict-then-optimise distinction~\cite{elmachtoub2022spo} in a
traffic setting. A rule trained to name the winner is rewarded for being right
in contexts where the margin is negligible and punished only mildly for being
wrong in contexts where it is large --- exactly the wrong trade under an
\NumAsymRatio$\times$ asymmetry.

\subsection{The learned selector works, and it is not worth deploying}

Table~\ref{tab:ladder} and Figure~\ref{fig:ladder} give the ladder. Predicting
\emph{advantage} rather than identity reverses the failure: B4 reaches
\NumBFourRegret~s of mean regret, closing \NumBFourGapClosed{} of the gap
between the fixed policy and the retrospective best, and it does so while being
\emph{less} accurate at naming the winner (\NumBFourTopOne{}) than the logistic
baseline (\NumBThreeTopOne{}, \NumBThreeRegret~s). It also has the lowest
maximum regret of any selector, \NumBFourMax~s against \NumBZeroMax~s.

The primary comparison declared in advance was B4 against B1. B4 wins it
decisively: \NumBFourRegret~s against \NumBOneRegret~s.

The magnitudes are what matter operationally. B4 improves mean system journey
time from \NumBZeroRegret~s of regret to \NumBFourRegret~s --- a saving of
\NumBFourVsBZero~s per vehicle on a mean journey time of
\NumFixedCostPThree~s. \textbf{That is under a tenth of one per cent.} The
selector is real, reproducible, and below the threshold at which any operator
would act on it.

\subsection{Abstention prevents every out-of-distribution harm, and forgoes the
in-distribution gain}

The selective selector B5 abstains in \NumAbstainRate{} of contexts and returns
exactly the fixed policy's regret, \NumBFiveRegret~s. The reason is visible in
the mechanism rather than in a tuning choice: the conformal interval on the
predicted advantage is wider than the advantage itself, so the requirement that
the lower bound exceed zero is never met. \textbf{A calibrated gate, correctly
applied, refuses to act on an effect smaller than its own uncertainty.}

Table~\ref{tab:ood} and Figure~\ref{fig:ood} show why this matters. On the
interpolation control (O3) the learned selector helps, cutting regret from
\NumOODThreeMidBZero~s to \NumOODThreeMidBFour~s. Under extrapolation and shift
it hurts: on unseen high demand it raises regret from
\NumOODOneHighBZero~s to \NumOODOneHighBFour~s, and under the concept shift of
an unseen disruption regime from \NumOODSevenIncBZero~s to
\NumOODSevenIncBFour~s. In every one of these cases B5 recovers the fixed
policy's performance exactly.

The support gate behaves as designed: it flags \NumOODFourLagSupport{},
\NumOODFivePenSupport{} and \NumOODSixGreenSupport{} of held-out contexts as
outside the training envelope on the lag, penetration and green-ratio
extrapolations respectively. The mechanistic rule degrades worst of all under
shift, reaching \NumOODFivePenBOne~s on unseen penetration.

One shift is instructive precisely because the gates do \emph{not} fire. When
the disruption is moved to a corridor that carries none in any training context
(O8), the support gate flags \NumOODEightNorthSupport{} of held-out contexts,
because the descriptors of those contexts are identical to descriptors seen in
training --- the incident's \emph{location} is not among them. The learned
selector happens to do well there (\NumOODEightNorthBFour~s against
\NumOODEightNorthBZero~s for the fixed policy), but it does so unguarded. A
support gate defined on a feature space cannot detect a shift that leaves that
feature space unchanged, and this is a concrete instance rather than a
theoretical caveat.

\subsection{The switching boundaries cannot be narrowed at this seed count}

Where the main grid showed the preferred policy changing between adjacent
demand levels, we ran three intermediate demand levels in each of the
\NumBoundN{} widest-margin cells.

\NumBoundMono{} of the \NumBoundN{} refined sequences are monotone, with a
single change point, and in the point estimate they narrow the switching
bracket from the \NumBoundGrid~veh/h grid spacing to a median of
\NumBoundWidth~veh/h. All \NumBoundPOneN{} of the shortest-to-load-balancing
transitions are monotone. Their change points sit in a
\NumBoundWidth~veh/h bracket whose \emph{position} differs by cell, from
\NumBoundPOneEarliestLo--\NumBoundPOneEarliestHi~veh/h in the earliest cell to
\NumBoundPOneLatestLo--\NumBoundPOneLatestHi~veh/h in the latest, with a median
of \NumBoundPOneMedLo--\NumBoundPOneMedHi~veh/h. That the position moves with
the cell is the point: the boundary is a property of the operating conditions,
not a single number for the network.

\textbf{None of these refined steps is seed-resolved.} The point estimates move
consistently, but at five seeds no individual step between adjacent refined
levels exceeds twice its own standard error. The remaining \NumBoundNonMono{}
sequences are non-monotone --- the winner changes back and forth across the
refined grid --- and \NumBoundNonMonoPFour{} of those involve the
reliability-aware policy.

We therefore report the transition as a bracket and state that it is a point
estimate, not a resolved boundary. The refinement tells us the shortest-path to
load-balancing reversal is a genuine, ordered regime change, and that the
apparent reliability-aware transitions near it are not distinguishable from
noise.

\subsection{The criteria disagree, and the policies specialise}

Over all \NumContextsAnalysed{} contexts, \NumParetoNonSingleton{} have a
non-singleton Pareto set across the three criteria. Non-domination is not
concentrated in one policy: P3 is non-dominated in \NumParetoPctPThree{} of
contexts, P4 in \NumParetoPctPFour{}, P2 in \NumParetoPctPTwo{} and P1 in
\NumParetoPctPOne{}.

The criteria are not interchangeable. Journey time and \CO{} correlate at
\NumCorrCOneCTwo{} and rank the four policies identically in
\NumSameOrderCOneCTwo{} of contexts; journey time and stopped delay correlate at
\NumCorrCOneCThree{} and agree on the ranking in only \NumSameOrderCOneCThree{}.
The best policy on \CO{} differs from the best on journey time in
\NumBestCTwoDiffersPct{} of contexts; on stopped delay it differs in
\NumBestCThreeDiffersPct{}.

The policies specialise by criterion rather than by regime. \NumSpecCOnePolicy{}
minimises journey time in \NumSpecCOneN{} contexts and \NumSpecCTwoPolicy{}
minimises \CO{} in \NumSpecCTwoN{}, but \textbf{\NumSpecCThreePolicy{}
minimises stopped delay in \NumSpecCThreeN{} contexts (\NumSpecCThreePct{})}
while minimising journey time in only \NumSpecCOnePFour{}. The
reliability-aware policy routes towards the steady, unsignalised corridor: its
vehicles spend less time stationary and more time moving.

\textbf{This is where the portfolio's complementarity actually lives.} On the
scalar decision cost the portfolio is near-degenerate; on the criterion vector
it is not.

\paragraph{Eco-routing remains excluded, on evidence.} \CO{} and journey time
correlate at \NumCorrCOneCTwo{} and select the same policy in
\NumSameOrderCOneCTwo{} of contexts. An objective nearly monotone in travel time
adds a label rather than a policy, and the numbers supporting that judgement are
reported here so a reader may disagree with the threshold rather than with an
assertion.

\subsection{Sensitivity}

Across the three declared preference profiles the preferred policy changes in
\NumPrefChanges{} of contexts. The journey-time choice is reproduced by the
time-weighted profile in \NumPrefAgreeTime{} of contexts, by the balanced
profile in \NumPrefAgreeBalanced{} and by the environment-weighted profile in
\NumPrefAgreeEnvironment{}. The decision is therefore stable under reasonable
preference assumptions, and where it is not, stopped delay is what moves it.

The two frozen policy parameters behave differently. The dispersion weight of
the reliability-aware policy is close to immaterial, changing that policy's own
mean journey time by at most \NumLambdaMaxChangePct{}; across all four
parameter changes the identity of the best policy is unchanged in
\NumWinnerStableMin{} to \NumWinnerStableMax{} of the contexts tested. The admissibility tolerance
of the load-balancing policy is not: loosening it from \NumEpsFrozen{} to
\NumEpsAlt{} improves that policy's own mean journey time by
\NumEpsGainPct{} (\NumEpsGainSec~s). The frozen value therefore
\emph{understates} how dominant the load-balancing policy is, which strengthens
rather than weakens the central conclusion. We report this rather than
retrofitting the better value.
